\documentclass[12pt,a4paper]{article}

\usepackage{fontawesome5}
\providecommand{\faPencilAlt}{\faIcon{pencil-alt}}

\usepackage{styles/couleurs}
\usepackage{styles/mise_en_page}
\usepackage{styles/boites}
\usepackage{styles/maths}
\usepackage{styles/figures}

\usepackage{enumitem}
\usepackage{array}
\usepackage{tabularx}
\usepackage{amsmath}
\usepackage{amssymb}
\usepackage{tikz}
\usepackage{xparse}

\usetikzlibrary{arrows.meta,calc}

\setlist[itemize]{leftmargin=1.4em,itemsep=0.15em,topsep=0.2em}
\setlist[enumerate]{leftmargin=*,itemsep=0.25em,topsep=0.25em}

\newcommand{\lignepointillee}{%
  \noindent\dotfill\par
  \vspace{0.45em}
}

\newcommand{\lignesreponse}[1]{%
  \begin{tcolorbox}[
    enhanced,
    colback=white,
    colframe=NeutreC!40,
    boxrule=0.5pt,
    arc=1.5mm,
    left=3mm,
    right=3mm,
    top=2mm,
    bottom=1mm,
    before skip=2pt,
    after skip=2pt
  ]
    \foreach \i in {1,...,#1}{%
      \lignepointillee
    }
  \end{tcolorbox}
}

\newcommand{\titreseq}[1]{%
\begin{center}
  {\sffamily\LARGE\bfseries\color{TitreC}Séquence 3 -- Nombres relatifs}\\[0.2em]
  {\sffamily\large\bfseries\color{SavoirC}#1}
\end{center}
}

\begin{document}

\titreseq{Exercices hors classe}

\remarque{Fiche donnée à l'avance. Ne faire que les exercices indiqués par le professeur.}

\section*{A. Déterminer un signe}

Indique seulement le signe du résultat.

\[
(-8)\times3
\qquad
(-4)\times(-7)
\qquad
18\div(-6)
\qquad
(-42)\div(-7)
\]

\[
(-2)\times3\times(-5)
\qquad
(-1)\times(-4)\times(-7)
\qquad
(-8)\times0\times(-21)
\]

\lignesreponse{5}

\section*{B. Produits et quotients simples}

Calcule.

\[
(-8)\times3
\qquad
(-4)\times(-7)
\qquad
(-2{,}5)\times4
\qquad
(-0{,}4)\times(-30)
\]

\[
18\div(-6)
\qquad
(-42)\div(-7)
\qquad
(-12{,}8)\div2
\qquad
7{,}2\div(-5)
\]

\lignesreponse{8}

\section*{C. Produits de plusieurs facteurs}

Calcule.

\[
(-3)\times(-2)\times(-2)\times2
\]

\[
3\times(-2)\times2\times(-2)
\]

\[
(-2)+(-4)\times(-5)
\]

\[
2-4\times(-2)\times(-2)
\]

\lignesreponse{7}

\section*{D. Enchaînements}

Calcule en détaillant.

\[
7+(-3)\times5
\]

\[
(-24)\div(-6)-3\times(-2)
\]

\[
\bigl(8-3\times(-2)\bigr)\div(-7)
\]

\lignesreponse{8}

\section*{E. Petit problème}

Une formule donne :

\[
t_F=1{,}8\times t_C+32.
\]

Calcule la température en degrés Fahrenheit correspondant à \(-5^\circ\mathrm{C}\).

Puis retrouve la température en degrés Celsius correspondant à \(41^\circ\mathrm{F}\).

\lignesreponse{7}

\end{document}
